\chapter{Couche 2 — Les données}

L'application a des comptes et des sessions, mais elle oublie tout au
redémarrage. Cinq briques pour qu'elle se souvienne --- et pour que ce dont elle
se souvient reste secret.

% -----------------------------------------------------------------------------
\outil{SQLAlchemy 2.0}{brique 10 · l'ORM}

\quoi
Un \emph{Object-Relational Mapper} : il fait correspondre une classe Python à
une table SQL. On écrit des objets, il écrit le SQL.

\pourquoi
L'alternative est d'écrire le SQL à la main. Trois raisons l'ont écartée.

\textbf{Les injections SQL.} Écrire
\cmd{f"SELECT * FROM users WHERE nom='\{saisie\}'"} laisse un utilisateur nommé
\cmd{' OR '1'='1} lire toute la table. L'ORM passe toujours par des paramètres
liés : l'injection devient structurellement impossible, pas seulement évitée par
discipline.

\textbf{Deux moteurs.} Le projet tourne sur SQLite en développement et sur
PostgreSQL en production. Leur SQL diffère sur les dates, les types, les
clauses. L'ORM écrit le dialecte qui convient.

\textbf{La vérification.} Une faute de frappe dans un nom de colonne devient une
erreur Python immédiate, au lieu d'une erreur SQL découverte en production.

On aurait pu prendre l'ORM de Django, mais il impose Django. Celui de SQLAlchemy
est le standard du monde Python hors Django, et sa version 2.0 apporte une
écriture typée que les analyseurs comprennent.

\role
\fichier{2\_webapp/models.py} décrit les sept tables. Chaque colonne est annotée
pour que \cmd{ruff} et l'éditeur sachent son type :

\begin{lstlisting}[language=Python]
class Rental(Base):
    __tablename__ = "rentals"
    id: Mapped[int] = mapped_column(Integer, primary_key=True)
    user_id: Mapped[int] = mapped_column(ForeignKey("users.id"), nullable=False)
    end_time: Mapped[datetime] = mapped_column(DateTime, nullable=False)
    __table_args__ = (Index("idx_rentals_end", "status", "end_time"),)
\end{lstlisting}

L'index de la dernière ligne n'est pas décoratif : c'est exactement la requête
que le Faucheur exécute toutes les dix secondes.

\apport
La plateforme se souvient. Et une famille entière d'attaques --- l'injection
SQL --- disparaît du périmètre.

% -----------------------------------------------------------------------------
\outil{Alembic}{brique 11 · les migrations}

\quoi
L'outil qui fait évoluer le schéma d'une base \emph{sans perdre les données}.
Chaque changement devient un fichier numéroté qui connaît son prédécesseur ; la
base retient où elle en est.

\pourquoi
Tant qu'on développe, on efface la base et on recommence. Dès qu'il y a des
données réelles --- des comptes, des locations en cours --- c'est impossible.

Sans Alembic, la seule option serait d'écrire les \cmd{ALTER TABLE} à la main et
de se souvenir lesquels ont été appliqués où. C'est exactement le genre de
mémoire humaine qui échoue au pire moment.

\role
\fichier{2\_webapp/migrations/} contient deux révisions : le schéma initial, puis
l'ajout de la table \cmd{archives}. La chaîne d'intégration exécute
\cmd{alembic check}, qui échoue si les modèles et les migrations divergent ---
impossible d'oublier de générer une migration.

\begin{piege}{SQLite ne sait pas modifier une colonne}
PostgreSQL accepte \cmd{ALTER TABLE ... ALTER COLUMN}. SQLite, non : il faut
créer une table neuve, recopier, supprimer, renommer. Alembic sait le faire,
mais seulement si on le lui demande par \cmd{render\_as\_batch=True}. Sans cette
option, les migrations passent sur un moteur et échouent sur l'autre.
\end{piege}

\apport
Le schéma peut évoluer en production. C'est ce qui a permis d'ajouter la
sauvegarde du travail --- une table de plus --- sur une base contenant déjà des
comptes et un historique.

% -----------------------------------------------------------------------------
\outil{SQLite}{brique 12 · la base de développement}

\quoi
Une base de données contenue dans un seul fichier, sans serveur à lancer.

\pourquoi
Pour développer et pour les tests, démarrer un serveur PostgreSQL est une
friction : il faut l'installer, le configurer, le lancer. SQLite est déjà dans
Python.

Le projet l'utilise en mode \textbf{WAL} (\emph{write-ahead logging}), qui
permet à plusieurs processus de lire pendant qu'un autre écrit --- nécessaire
puisque l'application et le Faucheur touchent la même base.

\role
Base par défaut quand \cmd{DATABASE\_URL} n'est pas défini. Les 151 tests
l'utilisent : chacun part d'un fichier vierge dans un dossier temporaire, ce qui
garantit qu'aucun test n'hérite de l'état d'un autre.

\apport
Les tests tournent en vingt secondes sur n'importe quelle machine, sans service
à installer. C'est ce qui rend l'intégration continue légère.

% -----------------------------------------------------------------------------
\outil{PostgreSQL 16}{brique 13 · la base de production}

\quoi
Un serveur de base de données complet, conçu pour plusieurs clients simultanés.

\pourquoi
SQLite suffit à un processus, pas à quatre ouvriers Gunicorn plus le Faucheur
qui écrivent en même temps. Son verrouillage porte sur le fichier entier : sous
charge, les écritures s'attendent.

MySQL aurait fait l'affaire ; PostgreSQL a été préféré pour deux raisons
concrètes ici : les \textbf{index partiels} --- la clause \cmd{WHERE} sur un
index unique, qui garantit l'unicité des ports sans bloquer leur
réutilisation --- et le respect strict des types.

\role
Service \cmd{db} de la pile compose, image \cmd{postgres:16-alpine}. Son port
\textbf{n'est pas publié} : seuls les services du réseau interne y accèdent. Il
démarre avec \cmd{--data-checksums}, qui détecte une corruption silencieuse du
disque.

\apport
La plateforme supporte la concurrence réelle. Et l'index unique partiel devient
une garantie du moteur, pas une intention du code.

% -----------------------------------------------------------------------------
\outil{cryptography · Fernet}{brique 14 · le coffre}

\quoi
Fernet est une recette de chiffrement symétrique authentifié : elle combine AES
en mode CBC et un HMAC. On ne peut ni lire sans la clé, ni modifier sans que
cela se voie.

\pourquoi
Le problème est particulier : le mot de passe root d'une machine louée doit
pouvoir être \textbf{réaffiché}. On ne peut donc pas le hacher --- le hachage est
à sens unique. Il faut un chiffrement réversible.

Écrire soi-même du chiffrement est l'erreur classique : on choisit mal le mode,
on réutilise un vecteur d'initialisation, on oublie l'authentification. Fernet
n'a \textbf{aucun paramètre à choisir} : c'est précisément son intérêt pour qui
n'est pas cryptographe.

\begin{notion}{Hacher ou chiffrer : deux besoins, deux outils}
Le mot de passe d'un \textbf{compte} est \textbf{haché} : on ne doit jamais
pouvoir le retrouver, seulement vérifier qu'un essai donne la même empreinte.

Le mot de passe \textbf{root d'une machine louée} est \textbf{chiffré} : il doit
pouvoir être réaffiché à son propriétaire.

Hacher le second le rendrait inutilisable ; chiffrer le premier serait une
faute.
\end{notion}

\role
\fichier{2\_webapp/crypto.py} charge deux clés \textbf{distinctes} : celle qui
signe les sessions Flask, et celle qui chiffre les mots de passe root. Les
séparer garantit qu'un vol de la première ne donne pas accès aux machines.

Les clés viennent de l'environnement --- Ansible les injecte depuis son Vault ---
ou, à défaut, d'un fichier local en permissions \cmd{0600}.

\gros[0.86]{zoom-secret.png}{Le mot de passe root, affiché une seule fois. Après
cet instant il n'existe plus qu'en base, chiffré.}

\gros[0.86]{zoom-coffre.png}{Le coffre rouvert après re-saisie du mot de passe
de compte : les accès redeviennent visibles cinq minutes, puis se remasquent.}

\begin{piege}{Une session Flask est signée, pas chiffrée}
Notre première version déposait le mot de passe root dans le cookie de session,
le temps de l'afficher. Or le contenu d'un cookie de session Flask est encodé en
base64 : \textbf{n'importe qui peut le lire} sans la clé. La signature empêche
seulement de le modifier.

On ne stocke plus que l'identifiant de la machine ; le secret est relu en base
et déchiffré au moment de l'affichage, après vérification du propriétaire.
\end{piege}

\apport
Un secret qu'on peut revoir sans jamais le stocker en clair, et qui ne transite
plus par le navigateur.

% -----------------------------------------------------------------------------
\section*{Ce que la plateforme sait faire après la couche 2}

Elle retient les comptes, les locations et les machines ; elle sait faire
évoluer son schéma sans perdre de données ; et elle garde des secrets qu'elle
peut réafficher. Il lui manque l'essentiel : les machines à louer.
